\documentclass[aspectratio=169, table]{beamer}

\usetheme{Madrid}
\usecolortheme{default}
\definecolor{ucbblue}{RGB}{0, 51, 102}

\setbeamercolor{structure}{fg=ucbblue}
\setbeamercolor*{palette primary}{fg=white,bg=ucbblue}
\setbeamercolor*{palette secondary}{fg=white,bg=ucbblue!80}
\setbeamercolor*{palette tertiary}{fg=white,bg=ucbblue!90}
\setbeamercolor{title}{fg=white, bg=ucbblue}
\setbeamercolor{frametitle}{fg=white, bg=ucbblue}
\setbeamercolor{item}{fg=ucbblue}

\usepackage[T1]{fontenc}
\usepackage[utf8]{inputenc}
\usepackage{tgpagella}
\usefonttheme{serif}
\usepackage[spanish, es-noshorthands]{babel}
\usepackage{amsmath, amssymb, bm}
\usepackage{graphicx}
\usepackage[most]{tcolorbox}
\usepackage{booktabs}
\usepackage{tikz}
\usepackage{pgfplots}
\pgfplotsset{compat=1.18}
\usepackage[american]{circuitikz}
\usetikzlibrary{shapes.geometric, arrows.meta, positioning, calc}

\title[Potencia y MOSFET]{Transición a Potencia: Clases de Salida y Fundamentos de MOSFET}
\subtitle{IMT-221 Circuitos Electrónicos III — Semana 10, Sesión 3}
\author[B. Quiroga Turdera]{M.Sc. Ing. Bernardo Quiroga Turdera}
\institute[UCB Tarija]{Universidad Católica Boliviana ``San Pablo'' — Sede Tarija}
\date{Semana 10}

\begin{document}

\begin{frame}
    \titlepage
\end{frame}

\begin{frame}{Transición Curricular: De la Señal a la Actuación}
    \begin{center}
    \begin{tikzpicture}[node distance=2.2cm, auto, thick, >=Stealth]
        \node[draw, rounded corners, fill=gray!20, align=center] (prot) {Protección\\(U1)};
        \node[draw, rounded corners, fill=gray!20, right of=prot, xshift=0.5cm, align=center] (sense) {Sensing\\(U2)};
        \node[draw, rounded corners, fill=gray!20, right of=sense, xshift=0.5cm, align=center] (ctrl) {Control\\(U3)};
        \node[draw, rounded corners, fill=ucbblue!20, right of=ctrl, xshift=0.5cm, align=center] (pwr) {\textbf{Potencia}\\\textbf{(U3)}};
        \node[draw, rounded corners, fill=red!10, right of=pwr, xshift=0.5cm, align=center] (load) {Carga};
        
        \draw[->] (prot) -- (sense);
        \draw[->] (sense) -- (ctrl);
        \draw[->] (ctrl) -- (pwr);
        \draw[->] (pwr) -- (load);
        
        \draw[->, dashed] (load.south) -- ++(0,-0.5) -| node[pos=0.25, below] {Feedback (U4)} (ctrl.south);
    \end{tikzpicture}
    \end{center}
    \vspace{0.5cm}
    \begin{tcolorbox}[colback=ucbblue!5, colframe=ucbblue, title=La Pregunta Técnica Central]
        ¿Por qué la electrónica que mide una corriente no es necesariamente la misma electrónica que puede entregar potencia a una carga?
    \end{tcolorbox}
\end{frame}

\begin{frame}{Sensing $\neq$ Actuation (Ejemplo 1)}
    \begin{columns}
        \begin{column}{0.48\textwidth}
            \textbf{Etapa de Sensado (Ej. Shunt)} \\
            Procesa señales pequeñas para mantener linealidad y no perturbar el sistema:
            \begin{align*}
                V_s &= 0.2\,\text{V} \\
                I_s &= 2\,\text{mA} \\
                P_s &= V_s I_s \\
                P_s &= (0.2)(0.002) = \mathbf{0.4\,\text{mW}}
            \end{align*}
        \end{column}
        \begin{column}{0.48\textwidth}
            \textbf{Etapa de Carga (Ej. Motor/Actuador)} \\
            Demanda energía significativa para realizar trabajo físico:
            \begin{align*}
                V_L &= 12\,\text{V} \\
                I_L &= 2\,\text{A} \\
                P_L &= V_L I_L \\
                P_L &= (12)(2) = \mathbf{24\,\text{W}}
            \end{align*}
        \end{column}
    \end{columns}
    \vspace{0.5cm}
    \centering
    \Large $\frac{P_L}{P_s} = 60,000 \implies$ \textbf{Prioridades de diseño radicalmente distintas.} \\
    \normalsize Aparecen: Disipación, eficiencia, protección térmica.
\end{frame}

\begin{frame}{Puente Funcional: Clases Lineales (A, B, AB)}
    \begin{table}[]
        \centering
        \begin{tabular}{@{}lccc@{}}
        \toprule
        \textbf{Aspecto} & \textbf{Clase A} & \textbf{Clase B} & \textbf{Clase AB} \\ \midrule
        Conducción por BJT & Ciclo completo & $\sim$ Medio ciclo & $>$ Medio ciclo \\
        Corriente de reposo ($I_Q$) & Alta & Idealmente cero & Pequeña/Moderada \\
        Distorsión de cruce & No dominante & Relevante & Reducida \\
        Disipación en reposo & Alta ($P_Q = V_Q I_Q$) & Baja & Mayor que B \\
        \textbf{Rol Conceptual} & \textbf{Linealidad} & \textbf{Eficiencia} & \textbf{Compromiso} \\ \bottomrule
        \end{tabular}
    \end{table}
\end{frame}

\begin{frame}{Visualización: Crossover Distortion (Distorsión de Cruce)}
    \centering
    \begin{tikzpicture}
        \begin{axis}[
            width=12cm, height=5cm,
            axis lines=middle,
            xmin=0, xmax=10, ymin=-1.5, ymax=1.5,
            xlabel={$t$}, ylabel={$v_o(t)$},
            xtick=\empty, ytick=\empty,
            legend style={at={(0.5,-0.1)}, anchor=north, legend columns=3, draw=none}
        ]
        % Ideal (Entrada referenciada)
        \addplot[thick, dashed, gray, domain=0:10, samples=100] {sin(deg(x))} ;
        % Clase B (Dead zone)
        \addplot[thick, red, domain=0:10, samples=200] {
            sin(deg(x)) > 0.3 ? sin(deg(x))-0.3 : (sin(deg(x)) < -0.3 ? sin(deg(x))+0.3 : 0)
        };
        % Clase AB (Reducida)
        \addplot[thick, blue, domain=0:10, samples=200] {
            sin(deg(x)) > 0.05 ? sin(deg(x))-0.05 : (sin(deg(x)) < -0.05 ? sin(deg(x))+0.05 : 0)
        };
        
        \legend{Entrada Ideal, Clase B (Notch severo), Clase AB (Notch reducido)}
        \end{axis}
    \end{tikzpicture}
    
    \vspace{0.2cm}
    La Clase AB inyecta una \textbf{pequeña corriente de bias} para solapar la conducción de ambos dispositivos alrededor de $v_{in} = 0$.
\end{frame}

\begin{frame}{Del Push-Pull Lineal al Switching}
    En una etapa \textbf{lineal}, el dispositivo pasa periodos largos con $v > 0$ e $i > 0$ simultáneamente.
    \begin{equation*}
        p(t) = v_{\text{disp}}(t) \cdot i_{\text{disp}}(t) \gg 0
    \end{equation*}
    
    Para actuar sobre altas potencias sin quemar el dispositivo, operamos como \textbf{Interruptor (Switching)}:
    \begin{columns}
        \begin{column}{0.48\textwidth}
            \begin{block}{Estado OFF (Abierto)}
                $i_D \approx 0$ \\
                $v_{DS} = V_{BUS}$ \\
                $P_{OFF} \approx (V_{BUS})(0) = 0$
            \end{block}
        \end{column}
        \begin{column}{0.48\textwidth}
            \begin{block}{Estado ON (Cerrado)}
                $v_{DS} \approx \text{pequeño}$ \\
                $i_D = I_{CARGA}$ \\
                $P_{ON} \approx (\text{pequeño})(I_{CARGA}) \approx 0$
            \end{block}
        \end{column}
    \end{columns}
\end{frame}

\begin{frame}{NMOS de Enriquecimiento: Terminales y Control}
    \begin{columns}
        \begin{column}{0.45\textwidth}
            \centering
            \begin{circuitikz}[scale=1.3, transform shape]
                \draw (0,0) node[nmos](M){};
                \node[above] at (M.D) {Drain ($D$)};
                \node[below] at (M.S) {Source ($S$)};
                \node[left] at (M.G) {Gate ($G$)};
                
                \draw[->, red, thick] (-0.5, 0.2) to[bend right] node[left] {$V_{GS}$} (-0.2, -0.7);
                \draw[->, blue, thick] (0.5, 0.6) to[bend left] node[right] {$V_{DS}$} (0.5, -0.6);
            \end{circuitikz}
        \end{column}
        \begin{column}{0.5\textwidth}
            \textbf{Variable de Control Electrostática:}
            \begin{equation*}
                \boxed{V_{GS} = V_G - V_S}
            \end{equation*}
            \textit{Precaución:} Nunca hables de "voltaje de gate" sin referenciar a dónde está conectado el Source.
            
            \vspace{0.3cm}
            \begin{itemize}
                \item Gate está \textbf{aislada eléctricamente} en DC.
                \item No requiere corriente continua para mantener conducción (como el BJT).
                \item \textbf{PERO} presenta capacitancia: exige corriente dinámica $i_G = \frac{dq_G}{dt}$ durante la conmutación.
            \end{itemize}
        \end{column}
    \end{columns}
\end{frame}

\begin{frame}{Regiones de Operación (Modelo Ideal Long-Channel)}
    \begin{columns}
        \begin{column}{0.45\textwidth}
            \textbf{1. Cutoff (Corte)}
            \begin{equation*}
                V_{GS} < V_{TH} \implies I_D \approx 0
            \end{equation*}
            \vspace{0.1cm}
            
            \textbf{2. Ohmic / Triode / Resistive} \\
            Para encendido eficiente (Switch ON):
            \begin{align*}
                V_{GS} &> V_{TH} \\
                0 \le V_{DS} &< V_{GS} - V_{TH}
            \end{align*}
            Se modela como resistencia ($R_{DS(on)}$).
            
            \vspace{0.3cm}
            \textbf{3. Saturation}
            \begin{align*}
                V_{GS} &> V_{TH} \\
                V_{DS} &\ge V_{GS} - V_{TH}
            \end{align*}
        \end{column}
        \begin{column}{0.5\textwidth}
            \centering
            \begin{tikzpicture}
                \begin{axis}[
                    width=7cm, height=6cm,
                    xmin=0, xmax=5, ymin=0, ymax=6,
                    xlabel={$V_{DS}$}, ylabel={$I_D$},
                    xtick=\empty, ytick=\empty,
                    axis lines=left,
                    domain=0:5
                ]
                % Ohmic dividing line
                \addplot[dashed, gray, domain=0:3] {x^2};
                \node[gray, rotate=60] at (1.5, 2.5) {$V_{DS} = V_{GS}-V_{TH}$};
                
                % Curves
                \addplot[thick, blue, domain=0:2] {2*(x - 0.5*x^2)};
                \addplot[thick, blue, domain=2:5] {2};
                
                \addplot[thick, ucbblue, domain=0:3] {3*(x - 0.5*x^2)};
                \addplot[thick, ucbblue, domain=3:5] {4.5};
                
                \node[align=center] at (1, 4.5) {\textbf{Ohmic}\\\textbf{(Switch ON)}};
                \node[align=center] at (4, 2.25) {\textbf{Saturation}};
                \end{axis}
            \end{tikzpicture}
        \end{column}
    \end{columns}
    \vspace{0.2cm}
    \begin{alertblock}{Distinción Crítica}
        \centering $\text{MOSFET Saturation} \neq \text{BJT Saturation}$ \\
        \small Para baja pérdida como switch, operamos en la región \textbf{Ohmic/Triode}, no en saturación.
    \end{alertblock}
\end{frame}

\begin{frame}{La Trampa del Threshold ($V_{GS(th)}$)}
    \begin{tcolorbox}[colback=red!5, colframe=red!70!black, title=Error Común]
        "El datasheet dice $V_{GS(th)} = 2\,\text{V}$. Si aplico $3.3\,\text{V}$ desde mi microcontrolador, el MOSFET estará totalmente encendido." $\implies$ \textbf{FALSO.}
    \end{tcolorbox}
    
    $V_{GS(th)}$ se define para una corriente de drenaje minúscula (inicio de conducción). \textbf{No demuestra baja resistencia ($R_{DS(on)}$).}
    
    \vspace{0.2cm}
    {\small
    \textbf{Ejemplo Ilustrativo (Datasheet IRLZ44N, Infineon):}
    \begin{itemize}
        \item $V_{GS(th)}$ típico: $1\,\text{V}$ a $2\,\text{V}$
        \item $R_{DS(on)}$ garantizada: $22\,\text{m}\Omega$ (@ $10\,\text{V}$), $25\,\text{m}\Omega$ (@ $5\,\text{V}$), $35\,\text{m}\Omega$ (@ $4\,\text{V}$)
    \end{itemize}
    }
    
    \vspace{0.2cm}
    \textbf{Conclusión:} Para diseñar, busca las condiciones bajo las cuales se especifica $R_{DS(on)}$, NO $V_{GS(th)}$.
\end{frame}

\begin{frame}{Pérdida por Conducción ($P_{cond}$)}
    Estando plenamente encendido (Región Ohmica), el MOSFET se modela como $R_{DS(on)}$.
    \begin{align*}
        V_{DS(on)} &\approx I_D \cdot R_{DS(on)} \\
        P_{cond} &= V_{DS(on)} \cdot I_D \quad \implies \quad \boxed{P_{cond} \approx I_D^2 \cdot R_{DS(on)}}
    \end{align*}
    
    \textbf{Ejemplo de Diseño (Evaluando opciones para $I_D = 4\,\text{A}$):}
    \begin{itemize}
        \item Candidato A: $R_{DS(on)} = 25\,\text{m}\Omega \implies P_A = (4)^2(0.025) = \mathbf{0.40\,\text{W}}$
        \item Candidato B: $R_{DS(on)} = 60\,\text{m}\Omega \implies P_B = (4)^2(0.060) = \mathbf{0.96\,\text{W}}$
    \end{itemize}
    \textit{Menor $R_{DS(on)}$ reduce las pérdidas estáticas de conducción... pero no es el único factor en PWM.}
\end{frame}

\begin{frame}{Corriente Dinámica de Compuerta y Gate Charge ($Q_g$)}
    La compuerta requiere carga para establecer el canal.
    \begin{equation*}
        Q_g = \int i_G(t) \,dt \quad \implies \quad I_{G,avg} \approx \frac{Q_g}{\Delta t}
    \end{equation*}
    
    \begin{center}
    \begin{tikzpicture}[node distance=2.5cm, auto, thick, >=Stealth]
        \node[draw, rounded corners, fill=white] (qg) {$Q_g$ alta};
        \node[draw, rounded corners, fill=white, right of=qg, xshift=1cm] (ig) {Requiere más $I_G$};
        \node[draw, rounded corners, fill=white, right of=ig, xshift=1.5cm, align=center] (dt) {$\Delta t$ aumenta\\(si el driver es débil)};
        
        \draw[->] (qg) -- (ig);
        \draw[->] (ig) -- (dt);
    \end{tikzpicture}
    \end{center}
    
    Transiciones lentas aumentan el solapamiento entre $v_{DS}$ e $i_D$, incrementando las \textbf{pérdidas de conmutación}.
\end{frame}

\begin{frame}{Intuición Visual: Pérdida de Conmutación ($E_{sw}$)}
    Durante la transición ON $\to$ OFF o OFF $\to$ ON, ambos $v_{DS}$ e $i_D$ son no nulos simultáneamente.
    \begin{equation*}
        p(t) = v_{DS}(t) \cdot i_D(t) \quad \implies \quad E_{sw} = \int_{t_{sw}} p(t) \,dt
    \end{equation*}
    
    \vspace{-0.2cm}
    \begin{center}
    \begin{tikzpicture}
        \begin{axis}[
            width=10cm, height=4cm,
            xmin=0, xmax=4, ymin=0, ymax=3,
            axis lines=left,
            xlabel={tiempo $t$}, ylabel={Magnitud},
            xtick=\empty, ytick=\empty,
            legend style={at={(0.5,1.2)}, anchor=south, legend columns=3, draw=none}
        ]
        % v_DS (cae)
        \addplot[thick, blue, domain=0:4, samples=100] {2.5/(1+exp(5*(x-2))) + 0.1};
        % i_D (sube)
        \addplot[thick, red, domain=0:4, samples=100] {2.5/(1+exp(-5*(x-2))) + 0.1};
        % p(t) (pulso)
        \addplot[thick, black, fill=gray!30, domain=0:4, samples=100] { (2.5/(1+exp(5*(x-2)))) * (2.5/(1+exp(-5*(x-2)))) };
        
        \node[anchor=west] at (2.2, 1.5) {$E_{sw}$ (Área bajo la curva)};
        \legend{$v_{DS}$ (Tensión), $i_D$ (Corriente), $p(t)$ (Potencia instantánea)}
        \end{axis}
    \end{tikzpicture}
    \end{center}
    \textit{Un MOSFET con muy baja $R_{DS(on)}$ pero altísima $Q_g$ puede calentar más a altas frecuencias (PWM) debido a este efecto.}
\end{frame}

\begin{frame}{Lectura Rigurosa de Datasheet (Datasheet Literacy)}
    No basta con un solo número. Debes revisar:
    
    \vspace{-0.2cm}
    \begin{table}[]
        \centering
        \resizebox{\textwidth}{!}{
        \begin{tabular}{@{}ll@{}}
        \toprule
        \textbf{Parámetro} & \textbf{Lectura Correcta de Ingeniería} \\ \midrule
        $V_{DSS}$ & Límite absoluto de tensión (Absolute Maximum Rating). \\
        $I_D$ (Continuous) & Siempre condicionado térmica/pulsos. No es corriente "segura" automática. \\
        $V_{GS(th)}$ & Inicio de conducción. \textbf{No} usar para diseño de encendido ON. \\
        $R_{DS(on)}$ & Leer siempre condicionado al $V_{GS}$ y Temp de prueba. \\
        $Q_g$ & Carga de compuerta; determina requerimientos del Gate Driver. \\
        Thermal ($R_{\theta JC}, R_{\theta JA}$) & Puente para diseño de disipación térmica. \\
        SOA (Safe Operating Area) & Frontera gráfica combinada de V, I y tiempo (pulsos). \\ \bottomrule
        \end{tabular}
        }
    \end{table}
\end{frame}

\begin{frame}{Próximos Pasos en el Curso}
    \begin{center}
        \Large $\text{MOSFET Switch} \rightarrow \text{Gate Driver} \rightarrow \text{Carga Inductiva} \rightarrow \text{H-Bridge}$
    \end{center}
    
    \vspace{0.5cm}
    \textbf{Preguntas abiertas para la siguiente sesión:}
    \begin{itemize}
        \item ¿Puede mi microcontrolador cargar la Gate lo suficientemente rápido?
        \item ¿Qué cambia drásticamente cuando la carga es \textit{inductiva} (ej. motor DC) en vez de resistiva?
        \item ¿Cómo implementamos topologías \textit{High-Side}?
        \item ¿Qué es el \textit{Shoot-Through} y cómo previene la destrucción del puente?
    \end{itemize}
\end{frame}

\end{document}
